\documentclass[11pt,a4paper]{article}
\usepackage[T1]{fontenc}
\usepackage{lmodern}
\usepackage[margin=22mm,headheight=15pt]{geometry}
\usepackage{amsmath,amssymb,mathtools,microtype}
\usepackage{xcolor,enumitem,fancyhdr,booktabs}
\definecolor{examblue}{HTML}{17365D}
\definecolor{rulegray}{HTML}{D6DCE4}
\pagestyle{fancy}
\fancyhf{}
\lhead{\textsc{Machine Learning 1}}
\rhead{\textsc{Mock Midterm 2 --- A}}
\cfoot{\thepage}
\setlength{\parindent}{0pt}
\setlength{\parskip}{0.45em}
\setlist[itemize]{leftmargin=1.5em,itemsep=0.2em}
\newcommand{\R}{\mathbb R}
\newcommand{\E}{\mathbb E}
\DeclareMathOperator{\Cov}{Cov}
\DeclareMathOperator{\tr}{tr}
\newcommand{\question}[3]{\par\medskip\noindent\textbf{#1}\quad\textbf{[#2 p]}\quad #3\par}
\newcommand{\answerbox}[1]{\par\smallskip\noindent\fcolorbox{rulegray}{white}{\parbox[t][#1][t]{\dimexpr\linewidth-2\fboxsep-2\fboxrule\relax}{\vspace{1mm}}}\par}
\newcommand{\exercise}[3]{\section*{\color{examblue}Exercise #1: #2}\textbf{#3}\par}

\begin{document}
\begin{center}
  {\Huge\bfseries\color{examblue}Machine Learning 1}\par
  \vspace{0.4em}
  {\LARGE\bfseries Mock Homework-based Midterm 2 --- A}\par
  \vspace{1em}
  {\large Two hours \qquad 18 points \qquad Two exercises}
\end{center}
\vspace{1em}
\textbf{Name:} \rule{8cm}{0.4pt}\hfill\textbf{Date:} \rule{3cm}{0.4pt}

\section*{Instructions}
\begin{itemize}
  \item This is an unofficial, deliberately challenging practice paper, not a prediction of the actual exam.
  \item Show your reasoning and intermediate steps. Partial credit is available. An unsupported final answer is not sufficient.
  \item For multivariate derivatives, first derive an arbitrary component in index notation, then assemble the vector or matrix.
  \item Use numerator layout: the Jacobian of $f:\R^m\to\R^n$ has shape $n\times m$. We write $\nabla_w f$ for the column gradient, so $df/dw=(\nabla_w f)^\top$ for scalar $f$.
  \item All vectors are columns. All logarithms are natural. Treat fitted parameters as fixed when making predictions.
  \item Use a basic calculator if needed. Exact expressions are acceptable. Do not use the answer key while attempting the paper.
  \item Suggested timing: 65 minutes for Exercise 1, 45 minutes for Exercise 2, and 10 minutes to check your work.
\end{itemize}

\begin{center}
\begin{tabular}{llr}
\toprule
Exercise & Topic & Points\\
\midrule
1 & Generative classification and regularized softmax & 10\\
2 & Affine PCA, a constrained component, and PPCA & 8\\
\bottomrule
\end{tabular}
\end{center}

\section*{Useful definitions}
\[
\sigma(a)=\frac{1}{1+e^{-a}},\qquad
\operatorname{softmax}(a)_k=\frac{e^{a_k}}{\sum_j e^{a_j}}.
\]
For symmetric positive definite $C\in\R^{D\times D}$,
\[
\mathcal N(x\mid\mu,C)=\frac{\exp[-\tfrac12(x-\mu)^\top C^{-1}(x-\mu)]}
{(2\pi)^{D/2}|C|^{1/2}}.
\]

\newpage
\exercise{1}{Two sensors, one class-dependent rate}{10 points}
A device belongs to one of $K$ classes. Two sensors give an input
$x=(m,s)^\top$, where $m\in\{0,1,\ldots\}$ is a count and $s\geq0$ is a waiting time.
The sensors are independent \emph{conditional on the class}. For class $C_k$,
\[
p(m\mid C_k)=\frac{(2\lambda_k)^m e^{-2\lambda_k}}{m!},\qquad
p(s\mid C_k)=\lambda_k e^{-\lambda_k s},\qquad \lambda_k>0.
\]
The same $\lambda_k$ occurs in both sensors; it is not two separate parameters.
The class priors $\pi_k>0$, with $\sum_k\pi_k=1$, are also unknown.
The training observations $(m_n,s_n,t_n)$ are independent, and $t_n$ is a
one-hot class vector. Define
\[
N_k=\sum_n t_{nk},\qquad M_k=\sum_n t_{nk}m_n,\qquad
S_k=\sum_n t_{nk}s_n.
\]
Assume $N_k>0$ for every class.

\question{1a}{3}{Write the joint likelihood of the observed inputs and labels, then its log-likelihood. Derive the maximum-likelihood estimates of \emph{both} $\pi_k$ and $\lambda_k$. Enforce the prior-probability constraint with a Lagrange multiplier and verify that the rate estimate is a maximum.}
\answerbox{10.0cm}

\newpage
\textbf{Exercise 1 continued}

For the remaining generative-model questions, use this training summary:
\begin{center}
\begin{tabular}{crrr}
\toprule
Class & $N_k$ & $M_k$ & $S_k$\\
\midrule
$C_1$ & 4 & 12 & 8\\
$C_2$ & 2 & 14 & 4\\
$C_3$ & 2 & 2 & 4\\
\bottomrule
\end{tabular}
\end{center}

\question{1b}{2}{Calculate the fitted priors and rates. Express the posterior probabilities as a softmax of scores affine in $(m,s)$. Give the two explicit inequalities that define the region in which $C_1$ has posterior probability at least as large as each competitor. If $m$ is temporarily allowed to be a nonnegative real number, is this region convex? Justify your answer.}
\answerbox{7.0cm}

\question{1c}{1}{For the complete input $(m,s)=(0,\tfrac12)$, identify the highest-posterior class. Now suppose the time sensor fails independently of the input and class, so only $m=0$ is observed. Derive the correct posterior from the count alone and identify its highest-posterior class. Explain why setting the missing waiting time to zero is not marginalization.}
\answerbox{5.0cm}

\newpage
\textbf{Exercise 1 continued}

\question{1d}{1}{Generate a new device by sampling its class from the fitted priors, then sample its two sensors from that class. You may use $\E[M\mid C_k]=2\lambda_k$ and $\E[S\mid C_k]=1/\lambda_k$. Derive and calculate the \emph{unconditional} covariance $\Cov(M,S)$. Does conditional independence imply unconditional independence here?}
\answerbox{4.0cm}

Next, fit a \emph{discriminative} classifier to the original labeled data, using
$\phi_n=(1,m_n,s_n)^\top$ and freely adjustable $w_k\in\R^3$:
\[
y_{nk}=\frac{\exp(w_k^\top\phi_n)}{\sum_j\exp(w_j^\top\phi_n)},\qquad
p(w_1,\ldots,w_K)=\prod_k\mathcal N(w_k\mid0,(\alpha R)^{-1}).
\]
Here $\alpha>0$ and $R\in\R^{3\times3}$ is a known symmetric positive definite
matrix. The prior includes the intercepts.

\question{1e}{2}{Obtain a negative log-posterior objective, omitting constants independent of the weights. Derive $\partial L/\partial w_{jr}$ in index notation and assemble $\nabla_{w_j}L$. Do not assume that $R$ is diagonal.}
\answerbox{5.0cm}

\question{1f}{1}{Without the prior, show that replacing every $w_k$ by $w_k+v$, for the same $v\in\R^3$, leaves the predictions unchanged. Prove that the MAP minimizer under the stated prior satisfies $\sum_k w_k=0$. Explain what the prior does to this parameter ambiguity.}
\answerbox{3.0cm}

\newpage
\exercise{2}{PCA with an offset and a forbidden direction}{8 points}
A dataset contains $J$ distinct vectors $x_i\in\R^3$. Vector $x_i$ occurs $n_i>0$
times among the $N=\sum_i n_i$ observations. Set $q_i=n_i/N$, so $\sum_iq_i=1$.
The full dataset has mean and covariance
\[
\bar x=\sum_iq_ix_i=\begin{pmatrix}4\\-2\\1\end{pmatrix},\qquad
S=\sum_iq_i(x_i-\bar x)(x_i-\bar x)^\top
=\begin{pmatrix}5&2&0\\2&5&0\\0&0&1\end{pmatrix}.
\]
The covariance uses the $1/N$, not the $1/(N-1)$, convention.
A one-dimensional \emph{affine} reconstruction uses a unit vector $u$ and an offset $c$:
\[
\widehat x_i=c+uu^\top(x_i-c),\qquad
E(c,u)=\sum_iq_i\|x_i-\widehat x_i\|^2,\qquad u^\top u=1.
\]

\question{2a}{2}{For a fixed $u$, derive \emph{all} minimizing offsets $c$. Show how the remaining minimization reduces to a trace expression involving $S$ and $uu^\top$. Explain why taking $c=\bar x$ is valid but does not make the offset unique.}
\answerbox{6.0cm}

\question{2b}{1}{With $c=\bar x$, find an optimal unconstrained unit vector $u$, the retained variance, and the minimum average reconstruction error. For the new input $x_*=(6,-1,3)^\top$, give its scalar coordinate $z_*=u^\top(x_*-\bar x)$, its reconstruction in $\R^3$, and its squared reconstruction error.}
\answerbox{5.0cm}

\newpage
\textbf{Exercise 2 continued}

\question{2c}{2}{A measurement policy now forbids the direction $v=(1,1,0)^\top/\sqrt2$: require $u^\top v=0$ as well as $u^\top u=1$. Use a Lagrangian to find all stationary directions of the retained variance. Select the global maximum and give the resulting average reconstruction error. Also reconstruct $x_*$ and calculate its squared error under this constraint.}
\answerbox{7.0cm}

For the next questions, remove the direction restriction and fit standard PPCA to
the full dataset of $N$ observations:
\[
z\sim\mathcal N(0,1),\qquad
x=wz+\mu+\epsilon,\qquad \epsilon\sim\mathcal N(0,\tau I_3),
\]
with independent $z,\epsilon$ and $\tau>0$. Fix $\mu=\bar x$ and constrain
$w=\sqrt a\,u_1$, where $a\geq0$ and $u_1$ is the leading eigenvector of $S$.
The marginal covariance is $C=a u_1u_1^\top+\tau I_3$.

\question{2d}{2}{Write the marginal log-likelihood as a function of $a$ and $\tau$, up to parameter-independent constants, using the eigenvalues of $S$. Derive its maximizing values of $a$ and $\tau$; do not only quote the PPCA formula. Give the fitted $C$ explicitly. Does it equal $S$? Explain the discrepancy in the discarded directions.}
\answerbox{6.0cm}

\newpage
\textbf{Exercise 2 continued}

\question{2e}{1}{An alternative parameterization uses $z'\sim\mathcal N(2,9)$ and $x=w'z'+\mu'+\epsilon$, with the \emph{same fitted noise variance} as in 2d. Find $w'$ and $\mu'$ that preserve the fitted marginal distribution of $x$. Could this new latent prior make a one-dimensional PPCA model match $S$ exactly? Justify your answer using covariance rank or eigenvalues.}
\answerbox{7.0cm}

\section*{Additional working space}
Label any answer continued here with its question number.
\answerbox{9.0cm}
\vfill
\begin{center}\textbf{End of Mock Midterm 2 --- A}\end{center}
\end{document}
